\section{生成式 EVLM 设计}

本章描述本文生成式 EVLM 的设计与实现. 在 A2 建立了 quantitative semantic prediction baseline 且 VTF 构建了 token-level EEG-enhanced vision tokens 之后, 本章将 EEG-enhanced vision tokens 输入预训练生成式 VLM, 实现从 EEG 信号到自由自然语言 caption 的完整生成通路.

\subsection{生成式 EVLM 的设计目标}

生成式 EVLM 的设计目标不是从 EEG 直接生成任意文本（``EEG-to-text''）, 而是: 给定一幅退化图像和对应的 EEG trial, 利用 EEG-enhanced vision tokens 作为视觉输入, 通过预训练生成式 VLM 生成关于图像内容的自然语言描述. 核心假设是: EEG-enhanced vision tokens 比纯视觉（vision-only）tokens 在退化条件下携带更准确的语义信息, 因此生成的 caption 应该具有更高的类别命中率.

\subsection{早期生成尝试与问题}

在项目初期（Day3 阶段）, 本文尝试了最简单直接的生成方案: 将 CLIP/EEG pooled embedding 通过 soft prompt projector（线性投影）映射为 $K=8$ 个软提示 token, 拼接在 Qwen2.5-1.5B-Instruct 的输入序列前部, 冻结语言模型仅训练投影器. 这种 soft prompt captioning 方案在 dummy data 上成功运行, 但当迁移到 real Thought2Text 数据时遇到了严重问题:

\begin{itemize}
  \item 模型生成大量 URL-like 文本（如 ``http://...''形式的字符串）;
  \item 生成 ImageNet wnid 代码（如 ``n04069434''）而非自然语言描述;
  \item 生成代码片段和不可读的 token 序列;
  \item 在退化条件下 caption 质量极低, 有效输出率不足 30\%.
\end{itemize}

这些问题的根源在于: (1) Qwen2.5 的训练数据中包含大量 URL 和代码, soft prompt 的训练信号不足以引导其生成自然的图像描述; (2) 仅 8 个软提示 token 的容量太小, 难以编码足够的视觉语义信息; (3) Qwen2.5 本身是纯文本 LLM, 未经过视觉-语言对齐训练, 不具备根据视觉特征生成描述的能力.

基于这些教训, 本文转向两条改进方向: (1) 使用原生视觉语言模型（Qwen2-VL）替代纯文本 LLM, 因为 Qwen2-VL 已经经过图文对齐预训练, 理解如何处理 visual tokens; (2) 增强视觉 tokens 的语义质量（通过 A2/VTF）, 而不是直接使用 CLIP pooled embedding.

\subsection{轻量 GRU Caption Decoder Baseline}

在正式引入 Qwen2-VL 之前, 本文先实现了一个轻量级的 GRU-based autoregressive caption decoder 作为生成式 baseline. 该 decoder 以 EEG-enhanced vision tokens 的平均池化向量为初始隐状态, 使用单层 GRU 逐步解码文本 token. Model 参数量约 2.8M, 在 BLIP pseudo-caption 上进行训练.

GRU baseline 能够生成 free-form caption, 其优势在于: 训练极快、不依赖大型预训练模型、输出始终为自然长度文本. 但 GRU decoder 不是预训练 VLM, 缺乏丰富的语言先验知识, 生成的 caption 在语法自然度和语义准确性上明显弱于 Qwen2-VL. GRU baseline 的主要价值在于验证 EEG-enhanced vision tokens 确实可以驱动 autoregressive caption 生成, 并作为后续 Qwen2-VL 结果的参照下界.

\subsection{预训练生成式模型路线探索}

在 GRU baseline 之上, 本文系统探索了五条生成式路线, 汇总于表~\ref{tab:route-comparison}.

\begin{table}[H]
  \centering
  \caption{生成式 EVLM 路线对比}
  \label{tab:route-comparison}
  \begin{tabularx}{\linewidth}{l l l X}
    \toprule
    路线 & 输入形式 & 桥接模块 & 结论 \\
    \midrule
    Route1 & enhanced tokens 池化后投影为 inputs\_embeds & 线性投影 + LoRA & 有效输出率低, 生成质量不稳定 \\
    Route2 & Q-Former/Perceiver bridge & Cross-attention bridge & 实现了但效果一般, 依赖更大的训练资源 \\
    Route3 & LLaVA-style mm\_projector & MLP 投影 & 工程兼容性较难, 被放弃 \\
    Route4 & BLIP-2/Q-Former style & Q-Former & 有一定效果, 但不如预期 \\
    Route5 & enhanced tokens + Qwen2-VL adapter & Prefix/Adapter + LoRA & 最终最佳生成路线 \\
    \bottomrule
  \end{tabularx}
\end{table}

Routes 1-4 的详细结果在第 8 章给出. 概括而言: Route1（Qwen2.5 inputs\_embeds + LoRA）有效输出率低; Route2（Q-Former bridge）需要更大的训练数据和显存, 在当前小样本条件下不如 Route5; Route3（LLaVA projector）与 Qwen2-VL 的兼容性存在工程问题; Route4（BLIP-2 style）在 clean 图像上表现尚可但退化条件下不稳定.

\subsection{Route5 Qwen2-VL 最终生成式架构}

Route5 是本文最终选定的生成式 EVLM 架构, 其结构如图~\ref{fig:route5} 所示.

\placeholderfigure{Route5 Qwen2-VL generative EVLM 架构图}{route5-qwenvl-placeholder}{5.0cm}

Route5 的核心设计如下:

\textbf{视觉输入}: 退化图像经 CLIP ViT 后得到 visual tokens $F_{\mathrm{vis}}\in\mathbb{R}^{B\times50\times512}$, 经过 VTF 与 EEG tokens 融合后得到 enhanced vision tokens $F_{\mathrm{enhanced}}\in\mathbb{R}^{B\times50\times512}$.

\textbf{桥接模块}: 增强视觉 tokens 通过 Qwen2-VL 内置的 vision-language adapter 进行维度对齐. Qwen2-VL 的 adapter 是一个预训练的 MLP projector, 将 ViT 输出映射到 LLM 的 embedding space.

\textbf{语义提示与退化类型信息}: 为了帮助模型进行受控生成, 在 input 中附加了 semantic top-$k$ prompt（由 A2 预测的最可能类别信息）和 corruption type label. 这些辅助信息以文本 token 形式插入生成 prompt 中.

\textbf{生成模型}: 使用 Qwen2-VL-2B-Instruct 作为生成 backbone. 该模型已经在大规模图文数据上进行过预训练, 具有视觉-语言理解和对齐能力. 模型参数约 2B, 单 GPU 可运行.

\textbf{LoRA 微调}: 在 Qwen2-VL 的注意力投影矩阵（Query 和 Value）上施加 LoRA, 秩 $r=8$. 可训练的 LoRA 参数数量仅约 1.3M, 远小于模型总参数量. 冻结所有原始模型权重, 仅训练 LoRA 适配器和 VTF cross-attention 模块.

\textbf{训练目标}: caption 训练使用 teacher-forcing 的自回归交叉熵损失. 目标序列为 ``corruption\_type: \{class\_name\}: caption'', 其中 caption 为 BLIP 伪 caption 或类别名（根据 target strategy）.

\textbf{推理}: 测试时, 以 ``corruption\_type:'' 为前缀触发生成, 模型自回归解码直到结束 token 或最大长度.

\subsection{LoRA 配置、Caption Target 消融与 Reranking}

\textbf{LoRA 秩选择}: 实验对比了 $r=8$ 和 $r=16$ 两种配置. 结果（第 8 章）表明 $r=8$ 在 valid caption rate 和 class-hit 上均优于 $r=16$. 推测原因是: 小样本条件下, 过大的可训练参数秩导致过拟合, $r=8$ 的低秩约束起到了隐式正则化的作用.

\textbf{Caption Target 策略}: 实验对比了两种训练目标配置:
\begin{itemize}
  \item \textbf{T1 (class-only)}: caption 仅包含类别名称, 目标序列为 ``\{class\_name\}''. 该策略直接以类别正确性为优化目标, 有利于 class-hit 提升;
  \item \textbf{T3 (class+BLIP)}: caption 包含类别名和 BLIP 生成的伪 caption, 目标序列为 ``\{class\_name\}: \{BLIP caption\}''. 该策略使生成文本更自然, 但 BLIP 伪 caption 中的噪声可能降低 class-hit.
\end{itemize}
T1 在 class-hit 上效果更好, T3 在生成文本的自然度上更优. 最终报告以 T3 为主展示生成效果.

\textbf{Best-of-N Reranking}: 推理时对每个输入生成 $N=5$ 个候选 caption（通过调整温度参数）, 从候选中选取 CLIP text prototype 分类命中最高的 caption 作为最终输出. Reranking 显著提升了 valid caption rate 和展示效果质量, 在 qualitative examples 中效果尤为明显.

\subsection{Route5 的定位}

需要再次强调: Route5 是本文构建的\textbf{生成式 EVLM 原型}, 不是自称的 ``EEG-to-text'' 系统. 它的主要贡献在于证明了: 在小样本 EEG-image 数据条件下, 将 EEG-enhanced vision tokens 注入预训练 Qwen2-VL 可以产生比纯视觉输入更准确的 free-form caption 生成. 然而, A2 在 quantitative classification 上的性能仍然远强于 Route5 的生成式结果, 说明在小样本条件下 constrained prediction 比 free-form generation 更可靠.
